# `main.tex` completo

```latex
\documentclass[12pt,a4paper]{article}

\usepackage[utf8]{inputenc}
\usepackage[T1]{fontenc}
\usepackage[spanish]{babel}

\usepackage[a4paper,margin=2.5cm]{geometry}

\usepackage{graphicx}
\usepackage{array}
\usepackage{longtable}
\usepackage{booktabs}
\usepackage{amsmath}
\usepackage{hyperref}
\usepackage{setspace}
\usepackage{float}
\usepackage{enumitem}

\onehalfspacing

\hypersetup{
    colorlinks=true,
    linkcolor=black,
    urlcolor=blue
}

\begin{document}

% ==============================
% PORTADA
% ==============================

\thispagestyle{empty}

\begin{center}

\vspace*{1cm}

{\Large\textbf{UNIVERSIDAD TÉCNICA ESTATAL DE QUEVEDO}\par}

\vspace{0.5cm}

{\large Facultad de Ciencias de la Ingeniería\par}
{\large Carrera de Ingeniería de Software\par}

\vspace{2cm}

{\Large\textbf{PRUEBA PRÁCTICA — UNIDAD IV}\par}

\vspace{0.5cm}

{\large Validación, Gestión, Herramientas y Estándares de Requisitos\par}

\vspace{2cm}

\textbf{Asignatura:} Ingeniería de Requisitos (ISR-401)

\vspace{0.5cm}

\textbf{Docente:} Ing. Gleiston Guerrero, Mg.

\vspace{0.5cm}

\textbf{Caso:} Sistema de Gestión de Pedidos

\vspace{1.5cm}

\textbf{Estudiante:} Robinson Espinoza Jeanpierre

\vspace{0.5cm}

\textbf{Fecha:} 11/08/2026

\vfill

\textbf{Link del Repositorio}

\vspace{0.3cm}

\url{https://github.com/jean200525/Evaluacion_UnidadIV.git}

\end{center}

\clearpage

% ==============================
% P1
% ==============================


\section{P1. Modelo de datos — Diagrama de clases UML}

El diagrama de clases representa el dominio del Sistema de Gestión de Pedidos mediante las clases Cliente, Pedido, LíneaPedido y Producto.

\begin{figure}[H]
    \centering
    \includegraphics[width=0.55\textwidth,height=0.45\textheight,keepaspectratio]{diagrama_clases.png}
    \caption{Diagrama de clases UML — Sistema de Gestión de Pedidos.}
    \label{fig:clases}
\end{figure}

% ==============================
% P2
% ==============================

\section{P2. Modelo funcional — Diagrama de actividades UML}

El proceso ``Registrar pedido'' se modela mediante un nodo inicial, una decisión relacionada con la disponibilidad de stock, dos caminos alternativos para confirmar el pedido o notificar el faltante, su convergencia y un nodo final. También se incluyen carriles de responsabilidad para Cliente y Sistema.

\begin{figure}[H]
    \centering
    \includegraphics[width=0.65\textwidth,height=0.45\textheight,keepaspectratio]{diagrama_actividades.png}
    \caption{Diagrama de actividades UML — Registrar pedido.}
    \label{fig:actividades}
\end{figure}

% ==============================
% P3
% ==============================

\section{P3. Modelo de comportamiento — Máquina de estados UML}

La máquina de estados representa el ciclo de vida de un Pedido mediante los estados Registrado, Confirmado, Enviado, Entregado y Cancelado.

Cada transición se encuentra etiquetada con su evento correspondiente. Además, se incluye la guarda \texttt{[stock disponible]} en \texttt{confirmar()} y la guarda \texttt{[no entregado]} en \texttt{anular()}. La transición de anulación se factoriza mediante el superestado EnProceso, que agrupa los estados Confirmado y Enviado.

\begin{figure}[H]
    \centering
    \includegraphics[width=0.65\textwidth,height=0.45\textheight,keepaspectratio]{maquina_estados.png}
    \caption{Máquina de estados UML — Ciclo de vida del Pedido.}
    \label{fig:estados}
\end{figure}

% ==============================
% P4
% ==============================

\section{P4. Consistencia entre las tres perspectivas}

La siguiente tabla presenta la verificación cruzada entre los eventos de la máquina de estados (P3), las funciones o acciones del diagrama de actividades (P2) y las entidades o atributos del diagrama de clases (P1).

\begin{longtable}{p{3cm}p{5cm}p{5cm}}
\toprule
\textbf{Evento (P3)} & \textbf{Función/acción (P2)} & \textbf{Entidad/atributo (P1)} \\
\midrule
\endfirsthead

\toprule
\textbf{Evento (P3)} & \textbf{Función/acción (P2)} & \textbf{Entidad/atributo (P1)} \\
\midrule
\endhead

\texttt{registrar()} &
Cliente ingresa datos del pedido &
Pedido (creación), Cliente (asociación) \\

\texttt{confirmar()} [stock disponible] &
Confirmar pedido &
Pedido.estado, Producto.stock (\texttt{descontarStock}) \\

--- (no existe en P3) &
Notificar faltante &
Producto.stock (\texttt{verificarStock}) \\

\texttt{despachar()} &
No modelado en P2 original &
Pedido.estado \\

\texttt{entregar()} &
No modelado en P2 original &
Pedido.estado \\

\texttt{anular()} &
No modelado en P2 original &
Pedido.estado \\

\bottomrule
\end{longtable}

\subsection*{Inconsistencia detectada}

El diagrama de actividades original (P2) solamente cubre el subproceso ``Registrar pedido'', hasta las acciones de confirmar o notificar faltante. Sin embargo, la máquina de estados (P3) incluye eventos posteriores como \texttt{despachar()}, \texttt{entregar()} y \texttt{anular()}, que no tenían una función equivalente modelada en el diagrama de actividades.

Por lo tanto, se identificó un defecto de tipo ``evento sin función''.

\subsection*{Corrección aplicada}

Se define un segundo proceso de actividades denominado ``Gestionar ciclo de vida del pedido'', en el cual se modelan explícitamente las actividades Despachar pedido, Entregar pedido y Anular pedido. Cada actividad dispara la transición correspondiente registrada en P3.

Con esta adición, los tres modelos quedan mutuamente consistentes.

% ==============================
% P5
% ==============================

\section{P5. Especificación de requisitos con esquema de atributos}

Se presentan dos requisitos funcionales y dos requisitos no funcionales. Cada requisito no funcional identifica una característica de ISO/IEC 25010:2023, establece una métrica y define un umbral.

\subsection{RF-01 — Requisito funcional}

El sistema debe permitir registrar un nuevo pedido asociando cliente, líneas de pedido y verificando disponibilidad de stock antes de confirmarlo.

\begin{longtable}{p{5cm}p{9cm}}
\toprule
\textbf{Atributo} & \textbf{Valor} \\
\midrule
Identificador & RF-01 \\
Prioridad & Alta \\
Estado & Propuesto \\
Fuente & Caso de negocio — módulo de gestión de pedidos \\
Estabilidad & Alta \\
Versión & 1.0 \\
Método de verificación & Prueba funcional (stock suficiente / insuficiente) \\
\bottomrule
\end{longtable}

\subsection{RF-02 — Requisito funcional}

El sistema debe permitir anular un pedido confirmado siempre que no haya sido entregado.

\begin{longtable}{p{5cm}p{9cm}}
\toprule
\textbf{Atributo} & \textbf{Valor} \\
\midrule
Identificador & RF-02 \\
Prioridad & Media \\
Estado & Propuesto \\
Fuente & Caso de negocio \\
Estabilidad & Media \\
Versión & 1.0 \\
Método de verificación & Prueba funcional (anular en distintos estados) \\
\bottomrule
\end{longtable}

\subsection{RNF-01 — Requisito no funcional — Eficiencia de desempeño}

El sistema debe registrar un pedido en menos de 3 segundos bajo carga normal, considerando hasta 50 usuarios concurrentes.

\begin{longtable}{p{5cm}p{9cm}}
\toprule
\textbf{Atributo} & \textbf{Valor} \\
\midrule
Característica & Eficiencia de desempeño (ISO/IEC 25010:2023) \\
Métrica & Tiempo de respuesta (segundos) \\
Umbral & $<3$ s \\
Identificador & RNF-01 \\
Prioridad & Alta \\
Estado & Propuesto \\
Fuente & Requisito de negocio (competitividad del servicio) \\
Estabilidad & Alta \\
Versión & 1.0 \\
Método de verificación & Prueba de rendimiento (carga) \\
\bottomrule
\end{longtable}

\subsection{RNF-02 — Requisito no funcional — Seguridad}

Los datos personales del cliente (cédula/RUC y correo) deben almacenarse cifrados y ser accesibles solamente mediante autenticación.

\begin{longtable}{p{5cm}p{9cm}}
\toprule
\textbf{Atributo} & \textbf{Valor} \\
\midrule
Característica & Seguridad (ISO/IEC 25010:2023) \\
Métrica & Porcentaje de campos sensibles cifrados en reposo \\
Umbral & 100\% de los campos sensibles cifrados \\
Identificador & RNF-02 \\
Prioridad & Alta \\
Estado & Propuesto \\
Fuente & Normativa de protección de datos \\
Estabilidad & Alta \\
Versión & 1.0 \\
Método de verificación & Auditoría de seguridad / prueba de penetración básica \\
\bottomrule
\end{longtable}

% ==============================
% P6
% ==============================

\section{P6. Priorización MoSCoW}

\begin{longtable}{p{3cm}p{3cm}p{9cm}}
\toprule
\textbf{Requisito} & \textbf{Categoría} & \textbf{Justificación} \\
\midrule
RF-01 & Must have &
Sin registrar pedidos con verificación de stock no existe el sistema; es el núcleo funcional. \\

RNF-02 & Must have &
Proteger los datos personales es obligatorio legal y contractualmente; no es negociable. \\

RF-02 & Should have &
Es importante para la operación y el control de errores/cancelaciones, pero el sistema puede operar sin ella en una primera versión. \\

RNF-01 & Could have &
Es deseable para la experiencia de usuario, pero no bloquea la operación básica si se entrega en una versión posterior. \\

\bottomrule
\end{longtable}

% ==============================
% P7
% ==============================

\section{P7. Validación por inspección}

Se aplica una lista de comprobación basada en ISO/IEC/IEEE 29148 para evaluar ambigüedad, completitud, singularidad, verificabilidad y factibilidad de los requisitos.

\subsection{Paso 1 — Identificación de defectos}

\begin{longtable}{p{1.5cm}p{1.5cm}p{1.5cm}p{1.5cm}p{1.5cm}p{1.5cm}p{5cm}}
\toprule
\textbf{Req.} &
\textbf{No ambiguo} &
\textbf{Completo} &
\textbf{Singular} &
\textbf{Verificable} &
\textbf{Factible} &
\textbf{Defecto detectado} \\
\midrule

RF-01 & Sí & No & Sí & Sí & Sí &
No especifica qué pasa si el cliente no existe (incompleto). \\

RF-02 & Sí & Sí & No & Sí & Sí &
Mezcla dos condiciones (confirmado / no entregado) en un solo enunciado; no es singular. \\

RNF-01 & Sí & Sí & Sí & Sí & No &
El umbral ``$<3$ s'' no aclara si incluye tiempo de red; factibilidad cuestionable. \\

RNF-02 & No & Sí & Sí & Sí & Sí &
``Cifrados'' es ambiguo: no especifica algoritmo o estándar mínimo. \\

\bottomrule
\end{longtable}

\subsection{Paso 2 — Retrabajo: requisitos corregidos}

\textbf{RF-01 (v1.1):} El sistema debe permitir registrar un pedido para un cliente previamente registrado, asociando una o más líneas de pedido, y validar la disponibilidad de stock de cada producto antes de confirmar el pedido.

\medskip

\textbf{RF-02a (v1.1):} El sistema debe permitir anular un pedido que se encuentre en estado Confirmado.

\medskip

\textbf{RF-02b (v1.1):} El sistema debe permitir anular un pedido que se encuentre en estado Enviado, siempre que no haya sido entregado.

\medskip

\textbf{RNF-01 (v1.1):} El sistema debe completar el registro de un pedido, desde el envío del formulario hasta la confirmación en pantalla, en menos de 3 segundos, medido en el servidor, bajo una carga de hasta 50 usuarios concurrentes.

\medskip

\textbf{RNF-02 (v1.1):} Los datos personales del cliente (cédula/RUC y correo) deben cifrarse en reposo usando AES-256 y ser accesibles solamente mediante autenticación válida.

% ==============================
% P8
% ==============================

\section{P8. Pruebas de aceptación trazadas}

\begin{longtable}{p{1.5cm}p{2cm}p{4cm}p{5cm}p{2cm}}
\toprule
\textbf{ID} & \textbf{Tipo} & \textbf{Entradas} & \textbf{Criterio de éxito} & \textbf{Traza a} \\
\midrule

PA-01 & Funcional &
Pedido con producto de stock=5, cantidad=2 &
El pedido queda en estado Confirmado y el stock se reduce a 3 &
RF-01 \\

PA-02 & Funcional &
Pedido en estado ``Confirmado'' &
Al anular, el pedido pasa a ``Cancelado'' y no puede volver a confirmarse &
RF-02 \\

PA-03 & No funcional &
50 usuarios concurrentes registrando pedidos &
El 95\% de las solicitudes se completa en $<3$ s &
RNF-01 \\

PA-04 & No funcional &
Consulta directa a BD del campo correo &
El campo se muestra cifrado, no legible en texto plano &
RNF-02 \\

\bottomrule
\end{longtable}

% ==============================
% P9
% ==============================

\section{P9. Matriz de trazabilidad}

\begin{longtable}{p{2cm}p{3cm}p{4cm}p{2.5cm}p{3cm}p{1.5cm}}
\toprule
\textbf{Requisito} &
\textbf{Fuente} &
\textbf{Modelo (traza adelante)} &
\textbf{Prueba} &
\textbf{Dependencia horizontal} &
\textbf{Huérfano} \\
\midrule

RF-01 &
Caso de negocio &
P1 (Pedido, LíneaPedido, Producto), P2 &
PA-01 &
Depende de RNF-02 &
No \\

RF-02 &
Caso de negocio &
P3 (máquina de estados) &
PA-02 &
-- &
No \\

RNF-01 &
Requisito de negocio &
P2 (actividad ``Registrar pedido'') &
PA-03 &
Relacionado con RF-01 &
No \\

RNF-02 &
Normativa protección de datos &
P1 (clase Cliente) &
PA-04 &
RF-01 depende de este &
No \\

\bottomrule
\end{longtable}

\subsection*{Resultado}

No se detectan requisitos huérfanos: los cuatro requisitos cuentan con al menos un modelo y una prueba de aceptación asociados.

% ==============================
% P10
% ==============================

\section{P10. Gestión del cambio y línea base}

\subsection{Solicitud de cambio SC-01}

\textbf{Descripción:} El área comercial solicita que RNF-01 reduzca el umbral de 3 segundos a 2 segundos para mejorar la competitividad del servicio.

\medskip

\textbf{Requisito afectado:} RNF-01.

\medskip

\textbf{Solicitante:} Área comercial.

\subsection{Análisis de impacto}

Al modificar RNF-01 se ve afectada la prueba PA-03, debido a que debe recalibrarse el umbral. También podría verse afectado RF-01, ya que ambos comparten el flujo ``Registrar pedido'' en P2.

No se afectan P1 ni P3.

\subsection{Decisión del CCB}

La solicitud es aprobada, condicionada a validar la factibilidad técnica de la infraestructura mediante una prueba de carga adicional antes de aceptar el nuevo umbral.

\subsection{Actualización de versión}

RNF-01 pasa de la versión 1.1 a la versión 1.2.

\subsection{Línea base (baseline) declarada}

La configuración congelada queda conformada por:

\begin{itemize}
    \item ERS v1.2 (con RNF-01 actualizado).
    \item Modelo de clases P1 v1.0.
    \item Modelo de actividades P2 v1.0.
    \item Máquina de estados P3 v1.0.
    \item Matriz de trazabilidad P9 v1.1.
    \item Conjunto de pruebas de aceptación v1.1 (PA-01 a PA-04, con PA-03 pendiente de recalibración posterior a la aprobación).
\end{itemize}

% ==============================
% REFERENCIAS
% ==============================

\section{Referencias}

\begin{enumerate}
    \item ISO/IEC/IEEE 29148, \textit{Systems and software engineering — Life cycle processes — Requirements engineering}.
    
    \item ISO/IEC 25010:2023, \textit{Systems and software engineering — Systems and software Quality Requirements and Evaluation (SQuaRE) — Product quality model}.
    
    \item OMG, \textit{Unified Modeling Language (UML)}.
\end{enumerate}

\end{document}
```
